\documentclass[11pt]{article}
\usepackage[margin=1in]{geometry}
\usepackage{amsmath,amssymb,amsthm,mathtools}
\usepackage{graphicx}
\usepackage{microtype}
\usepackage{needspace}
\usepackage[hidelinks]{hyperref}

\newtheorem{theorem}{Theorem}
\newtheorem{lemma}[theorem]{Lemma}
\newtheorem{remark}[theorem]{Remark}
\newcommand{\R}{\mathbb R}
\newcommand{\E}{\mathbb E}
\newcommand{\Id}{\operatorname{Id}}
\newcommand{\Hess}{D^2}
\newcommand{\comm}[2]{[#1,#2]}

\title{A Standard-Gaussian Counterexample to Equality of the\\
Kim--Milman and Brenier Transport Maps}
\author{Automated mathematics run \texttt{fa\_banach\_001}, lane 13\\
Model: GPT5.6}
\date{22 July 2026}

\begin{document}
\maketitle

\begin{center}
\fbox{\parbox{0.91\textwidth}{\textbf{Status.}
Candidate counterexample, likely valid, awaiting human review.  The proof is
analytic; the included symbolic computation only checks its elementary matrix
and integral arithmetic.}}
\end{center}

\section{Source question}

Kim and Milman constructed a contraction from a Gaussian measure to a more
log-concave measure by reversing an Ornstein--Uhlenbeck heat-flow
interpolation~\cite{KimMilman}.  Their Section~6 asked whether this map must
coincide with the Brenier optimal-transport map. Tanana proved that the answer
is no for a non-standard Gaussian source~\cite{Tanana}.  The final section of
her note observed that the two maps do coincide when both the source is
standard Gaussian and the target is Gaussian, but supplied numerical evidence
that they need not coincide for a general log-concave perturbation of the
standard Gaussian. The relevant page-5 passage is reproduced below.

\begin{figure}[ht]
  \centering
  \includegraphics[width=\textwidth]{figures/open_problem_crop.png}
  \caption{Tanana's standard-Gaussian question and numerical evidence
  \cite[p.~5]{Tanana}.}
\end{figure}

The question addressed here is therefore:

\begin{quote}
If the source is the standard Gaussian measure and
$d\nu=Z^{-1}e^{-F}d\gamma$ with $F$ convex, must the Kim--Milman map equal the
Brenier map?
\end{quote}

\section{Definitions}

Let $\gamma=\gamma_2$ be standard Gaussian measure on $\R^2$, and let
$(P_t)_{t\geq0}$ be its Ornstein--Uhlenbeck semigroup:
\[
 P_tg(x)=\int_{\R^2}g\bigl(e^{-t}x+\sqrt{1-e^{-2t}}\,y\bigr)\,d\gamma(y).
\]
Given a smooth convex globally Lipschitz $F$ and $\varepsilon>0$, put
\[
 d\nu_\varepsilon=Z_\varepsilon^{-1}e^{-\varepsilon F}\,d\gamma.
\]
Following Kim--Milman, define
\[
 Z_t^\varepsilon=-\log P_t(e^{-\varepsilon F}),\qquad
 \frac{d}{dt}S_t^\varepsilon(x)
   =\nabla Z_t^\varepsilon(S_t^\varepsilon(x)),\qquad
 S_0^\varepsilon=\Id.
\]
Then $S_t^\varepsilon$ transports $\nu_\varepsilon$ to the interpolating
measure and $S_\infty^\varepsilon$ transports $\nu_\varepsilon$ to $\gamma$.
The forward Kim--Milman map is
$T_{KM}^\varepsilon=(S_\infty^\varepsilon)^{-1}$.

\section{Proof intuition}

At first order in $\varepsilon$, the derivative of the inverse flow is an
integral of Hessians and is necessarily symmetric.  At second order, most
terms are still Hessians, but chronological multiplication remembers the
order in which different Hessians act.  Its skew part is exactly a double
integral of commutators
\[
 \int_{0<s<t<\infty}\comm{\Hess P_tF(x)}{\Hess P_sF(x)}\,ds\,dt.
\]
Thus it suffices to make two OU-smoothed Hessians fail to commute in an
integrated sense.  A sum of fourth powers of two nonorthogonal linear forms
does so explicitly.  We first calculate with that polynomial, then replace
the scalar fourth power by a smooth convex linear-tail truncation.  The
commutator coefficient varies continuously under this replacement, whereas
global Lipschitz regularity makes the Kim--Milman flow unproblematic.  A
nonsymmetric inverse-flow derivative is impossible for the inverse of a
Brenier map.

\Needspace{9\baselineskip}
\section{The perturbative skew formula}

\begin{lemma}[Well-posedness and integrable parameter expansion]\label{lem:bounds}
Let $F\in C^\infty(\R^2)$ be convex, bounded below, and globally Lipschitz,
with bounded derivatives of orders at least two.  For all sufficiently small
$\varepsilon\geq0$, the flow $S_t^\varepsilon$ and its spatial derivative are
defined for all $t\geq0$, converge as $t\to\infty$, and are twice
differentiable in $\varepsilon$ at $0$.  On every fixed compact set, all
remainders needed below have an integrable-in-$t$ envelope.
\end{lemma}

\begin{proof}
The OU commutation identity gives
\[
 \nabla P_tg=e^{-t}P_t(\nabla g),\qquad
 D^kP_tg=e^{-kt}P_t(D^kg).
\]
For $g=e^{-\varepsilon F}$,
\[
 \lvert\nabla Z_t^\varepsilon\rvert
 =\frac{\lvert\nabla P_tg\rvert}{P_tg}
 \leq \varepsilon e^{-t}\lVert\nabla F\rVert_\infty.
\]
The quotient rule, the same commutation identities, and boundedness of the
derivatives of $F$ give, locally uniformly in space and uniformly for small
$\varepsilon$,
\[
 \lVert D^kZ_t^\varepsilon\rVert\leq C_k\varepsilon e^{-kt}
 \quad(k=2,3),
\]
with harmless additional products having faster decay.  The denominators are
uniformly separated from zero on compact sets: after adding a constant assume
$F\geq0$, and use Jensen's inequality and the linear-growth bound on $F$ to
bound $P_t(e^{-\varepsilon F})$ below.

The velocity is integrable in time, while its derivative has an integrable
bound. Standard ODE estimates therefore give global existence at every fixed
initial point, convergence of $S_t^\varepsilon$ and $DS_t^\varepsilon$, and
uniform compact-set bounds. Differentiating
$-\log P_t(e^{-\varepsilon F})$ in $\varepsilon$ is justified by Gaussian
dominated convergence (a Lipschitz function has only linear growth). The same
integrable estimates permit twice differentiating the ODE and passing to
$t=\infty$. This proves all assertions.
\end{proof}

\begin{lemma}[Second-order skew formula]\label{lem:skew}
Under the hypotheses of Lemma~\ref{lem:bounds}, for every $x\in\R^2$,
\begin{equation}\label{eq:skew}
 DS_\infty^\varepsilon(x)-DS_\infty^\varepsilon(x)^T
 =\varepsilon^2\mathcal C_F(x)+o(\varepsilon^2),
\end{equation}
where
\begin{equation}\label{eq:C}
 \mathcal C_F(x)=
 \int_0^\infty\int_0^t
 \comm{\Hess P_tF(x)}{\Hess P_sF(x)}\,ds\,dt.
\end{equation}
\end{lemma}

\begin{proof}
Write $H_t=\Hess P_tF(x)$. Since
\[
 -\log P_t(e^{-\varepsilon F})
   =\varepsilon P_tF+O(\varepsilon^2),
\]
we have, along the perturbed trajectory,
\[
 \Hess Z_t^\varepsilon(S_t^\varepsilon(x))
   =\varepsilon H_t+\varepsilon^2K_t+o(\varepsilon^2).
\]
Here $K_t$ is symmetric: it is the sum of the Hessian of the direct second
$\varepsilon$-variation of $Z_t^\varepsilon$ and the directional derivative
of the Hessian $\Hess P_tF$ in the first variation of the trajectory.  The
latter remains symmetric in its two matrix indices because third derivatives
of a scalar function are symmetric.

Let $J_t^\varepsilon=DS_t^\varepsilon(x)$ and write
\[
 J_t^\varepsilon=I+\varepsilon A_t+\varepsilon^2B_t+o(\varepsilon^2).
\]
The variational equation
\[
 (J_t^\varepsilon)'=
 \Hess Z_t^\varepsilon(S_t^\varepsilon(x))J_t^\varepsilon
\]
gives
\[
 A_t=\int_0^tH_s\,ds,
 \qquad B_t'=K_t+H_tA_t.
\]
Both $K_t$ and $A_t$ are symmetric. Consequently
\[
 B_\infty-B_\infty^T
 =\int_0^\infty\left(H_tA_t-A_tH_t\right)dt
 =\int_0^\infty\int_0^t\comm{H_t}{H_s}\,ds\,dt.
\]
Lemma~\ref{lem:bounds} justifies the integrations and the limiting expansion,
which proves \eqref{eq:skew}.
\end{proof}

\section{An exact nonzero coefficient}

Let
\[
 a=\begin{pmatrix}1\\0\end{pmatrix},\qquad
 b=\begin{pmatrix}1\\1\end{pmatrix},\qquad
 F_0(x)=(a\cdot x)^4+(b\cdot x)^4,
 \qquad x_0=\begin{pmatrix}0\\1\end{pmatrix}.
\]
Set
\[
 P=\sum_{v\in\{a,b\}}(v\cdot x_0)^2vv^T=bb^T
 =\begin{pmatrix}1&1\\1&1\end{pmatrix},
\]
and
\[
 Q=\sum_{v\in\{a,b\}}\lvert v\rvert^2vv^T=aa^T+2bb^T
 =\begin{pmatrix}3&2\\2&2\end{pmatrix}.
\]

\begin{lemma}\label{lem:quartic}
For the polynomial model above,
\[
 \mathcal C_{F_0}(x_0)
 =\begin{pmatrix}0&6\\-6&0\end{pmatrix}\neq0.
\]
\end{lemma}

\begin{proof}
Put $r_t=e^{-2t}$.  The Mehler formula and the fourth Gaussian moment give,
for any fixed vector $v$,
\[
 D^2P_t(v\cdot x)^4
 =12\left(r_t^2(v\cdot x)^2+r_t(1-r_t)\lvert v\rvert^2\right)vv^T.
\]
Therefore
\[
 H_t:=\Hess P_tF_0(x_0)
 =12\left(r_t^2P+r_t(1-r_t)Q\right).
\]
It follows that
\[
 \comm{H_t}{H_s}
 =144r_tr_s(r_t-r_s)\comm{P}{Q}.
\]
Direct integration yields
\[
 \int_0^\infty\int_0^t
 r_tr_s(r_t-r_s)\,ds\,dt=-\frac1{24},
\]
whereas
\[
 \comm{P}{Q}=\begin{pmatrix}0&-1\\1&0\end{pmatrix}.
\]
Substitution in \eqref{eq:C} gives the claimed matrix.
\end{proof}

\section{Convex Lipschitz smoothing}

The polynomial $F_0$ is not globally Lipschitz, so we transfer the nonzero
coefficient to a Lipschitz convex model.

For $R>1$, define the even convex $C^1$ function
\[
 \widetilde\psi_R(u)=
 \begin{cases}
 u^4,& |u|\leq R,\\
 4R^3|u|-3R^4,& |u|\geq R.
 \end{cases}
\]
Convolve it with a nonnegative even $C^\infty$ mollifier of radius
$\delta_R\downarrow0$ and call the result $\psi_R$.  Then $\psi_R$ is smooth,
convex, globally Lipschitz, and has bounded derivatives of every positive
order. Define
\[
 F_R(x)=\psi_R(a\cdot x)+\psi_R(b\cdot x).
\]

\begin{lemma}\label{lem:smooth}
The mollification radii may be chosen so that
\[
 \mathcal C_{F_R}(x_0)\longrightarrow\mathcal C_{F_0}(x_0)
 \qquad(R\to\infty).
\]
In particular, $\mathcal C_{F_R}(x_0)\neq0$ for all sufficiently large $R$.
\end{lemma}

\begin{proof}
Choose $\delta_R\to0$ slowly enough that $\psi_R''(u)\to12u^2$ at every
fixed $u$.  The piecewise function has second derivative $12u^2$ inside
$(-R,R)$ and zero outside, apart from its corner in the distributional
description.  Its mollifications can be chosen with a uniform quadratic
envelope
\[
 0\leq\psi_R''(u)\leq C(1+u^2).
\]
For a twice differentiable function of at most polynomial growth,
\[
 \Hess P_tG(x)=e^{-2t}P_t(\Hess G)(x).
\]
Gaussian dominated convergence hence gives
\[
 \Hess P_tF_R(x_0)\longrightarrow\Hess P_tF_0(x_0)
 \quad\text{for every }t\geq0,
\]
and the quadratic envelope gives the uniform bound
\[
 \lVert\Hess P_tF_R(x_0)\rVert\leq Ce^{-2t}.
\]
Thus the norm of the commutator integrand is bounded by
$Ce^{-2t}e^{-2s}$ on $0<s<t<\infty$.  Dominated convergence in
\eqref{eq:C}, together with Lemma~\ref{lem:quartic}, proves the result.
\end{proof}

\section{Counterexample theorem}

\begin{theorem}\label{thm:main}
There exists a smooth convex globally Lipschitz function
$F:\R^2\to\R$ and $\varepsilon_0>0$ such that, for every
$0<\varepsilon<\varepsilon_0$, the Kim--Milman map transporting $\gamma_2$
to
\[
 d\nu_\varepsilon=Z_\varepsilon^{-1}e^{-\varepsilon F}\,d\gamma_2
\]
does not coincide with the Brenier map.
\end{theorem}

\begin{proof}
Choose $R$ as in Lemma~\ref{lem:smooth} and set $F=F_R$.  By
Lemma~\ref{lem:skew},
\[
 DS_\infty^\varepsilon(x_0)-DS_\infty^\varepsilon(x_0)^T
 =\varepsilon^2\mathcal C_F(x_0)+o(\varepsilon^2).
\]
Since $\mathcal C_F(x_0)\neq0$, the derivative
$DS_\infty^\varepsilon(x_0)$ is nonsymmetric for all sufficiently small
positive $\varepsilon$.

Suppose that the forward Kim--Milman map were the Brenier map.  The maps here
are smooth diffeomorphisms: this follows from the smooth globally Lipschitz
potential and the integrable bounds in Lemma~\ref{lem:bounds}.  The inverse of
a Brenier diffeomorphism $\nabla\phi$ is $\nabla\phi^*$, the gradient of the
convex conjugate. Its derivative is symmetric wherever defined.  This
contradicts the nonsymmetry of $DS_\infty^\varepsilon(x_0)$. Hence the two
forward maps are different.
\end{proof}

\section{Verification, limitations, and novelty}

The exact algebra in Lemma~\ref{lem:quartic} is independently checked by
\texttt{code/verify\_commutator.py} using symbolic rational arithmetic.  The
script is supplementary: the smoothing and perturbative ODE arguments are
analytic and do not depend on numerical evidence.

The theorem is existential: it proves that every sufficiently late smooth
linear-tail truncation supplies a counterexample for all sufficiently small
$\varepsilon$, but it does not provide optimized numerical values of $R$,
$\delta_R$, or $\varepsilon_0$.  It answers map equality only; both maps remain
valid contractions transporting the same measures.

The bounded novelty search on 22 July 2026 covered the run's cheap indexes and
arXiv/web searches for the exact ids 1002.0373 and 1709.06464 and the phrases
``Kim--Milman map standard Gaussian Brenier'', ``heat flow interpolation
Brenier standard Gaussian'', and close variants.  It found the source note
and later applications of the Kim--Milman construction, but no rigorous later
resolution of the standard-Gaussian/non-Gaussian question.  This is not an
exhaustive guarantee; novelty confidence is medium.

\paragraph{Human-review recommendation.}
First inspect the uniform-in-time parameter differentiation in
Lemma~\ref{lem:bounds}; then inspect the common quadratic envelope and
dominated-convergence passage in Lemma~\ref{lem:smooth}. The exact
commutator computation has low verification risk.

\begin{thebibliography}{9}
\bibitem{KimMilman}
Y.-H. Kim and E. Milman,
\emph{A Generalization of Caffarelli's Contraction Theorem via (reverse) Heat
Flow}, arXiv:1002.0373; Math. Ann. 354 (2012), 827--862.

\bibitem{Tanana}
A. Tanana,
\emph{Comparison of transport map generated by heat flow interpolation and
the optimal transport Brenier map}, arXiv:1709.06464 (2017).
\end{thebibliography}

\end{document}
